% ===== sec/11_conclusion.tex =====
\section{Conclusion}
\label{sec:conclusion}

Koohestani et al.~\cite{koohestani2025automata} asked whether
agents are just automata; we accept their answer and move on to
the operational question: can the Chomsky type of one agent be
used to bound the type of another?  We have given a self-contained
formal theory that answers \emph{yes}, derived from
Turing's~\cite{turing1936computable} 1936 theory of automata and
Chomsky's 1956--59 hierarchy~\cite{chomsky1956three,
chomsky1959certain}, in which an agent is identified with, and
only with, its trace language.  Within this theory:
\begin{itemize}
\item Agents are 7-tuples
$(\Sigma_\mathcal{A},Q_\mathcal{A},q_0^\mathcal{A},X_\mathcal{A},
\delta_\mathcal{A},F_\mathcal{A},h_\mathcal{A})$ over a fixed
operation alphabet
(Section~\ref{sec:trace}, Definition~\ref{def:agent}).
\item Trace languages classify into the standard
Chomsky hierarchy (Section~\ref{sec:chomsky}).  Constant-bounded
iteration loops are Type-3 (regular); they do not require indexed
grammars.
\item The producer--consumer construction of
Dijkstra~\cite{dijkstra1968cooperating}, instantiated as a
Type-0 Goal-Driven producer (idealised) bound by a Type-3 Data-Driven
pipeline DFA verifier, suffices: their intersection is closed at the
producer's class and decidable in linear time at the consumer's
DFA (Theorem~\ref{thm:closure}, Corollary~\ref{cor:budget}).
This closure theorem is a known result of automata theory
(Hopcroft--Ullman~\cite{hopcroft2006automata}); the contribution
of this paper is its application as an \emph{architectural
principle} for multi-agent systems.
\item The Type-3 consumer is realised as the deterministic
pipeline DFA implemented in \texttt{framework/neurogolf\_dfa\_verifier.py}
($|Q|=13$, $|\delta|=83$, decision time $0.684\,\mu$s per step).
\item Type-2 and Type-1 sub-agents are not designed; they are
projected from the binding (Theorem~\ref{thm:derived}).
\item A Type-3 agent may itself act as a producer,
backward-inferring constrained changes to the tape, provided the
inference terminates in a decidable verdict; otherwise the
halting problem returns and the verification budget collapses
(Section~\ref{sec:type3-producer}).
\item The Type-0 classification of the Goal-Driven agent is an
idealised upper bound; the physical machine is a linear-bounded
automaton, but every trace it produces is also producible by the
idealised Type-0 machine (Lemma~\ref{lem:bounded-soundness}).
\end{itemize}

The architecture is not metaphorical.  Instantiated on the
\textsc{NeuroGolf ONNX ARC-AGI} task, it produced a public
leaderboard score of $5793.14$.  Instantiated on the
\textsc{Spaceship Titanic} task from the public
\textsc{MLE-bench} benchmark, it produced a leaderboard placement
of $710$ out of $2330$ teams ($\approx\!30.5\%$, top tertile).
The Goal-Driven producer had no a priori knowledge of the task-specific
goal; the Data-Driven consumer had no LLM component.  The start
language has 7 sub-agents over 9 operation symbols
(\texttt{experiment\_6/trace\_language\_start.csv}); the realised
trace has 11 sub-agents over 17 operation symbols
(\texttt{experiment\_6/trace\_language.csv}), with 4 agent roles
and 9 operation symbols absent from the seed, all forced into
existence by the verifier's transition requirements.  Reproducible
code, the start-language CSV, the realised trace CSV, and the DFA
verifier are released with this paper.

\subsection{Future work}

Four directions are immediate.  First, the architecture should
be evaluated on a broader range of tasks to test the
external-validity claim of Section~\ref{sec:discussion}.  Second,
the formal Chomsky classification of the derived sub-agents should
be established via pumping-lemma proofs rather than structural
inspection.  Third, the trace-language identity criterion of
Definition~\ref{def:traceeq} invites a formal study of agent
equivalence: when are two agents with different LLM backbones
trace-equivalent on a given task?  Fourth, the characterisation
of regular transition-table update rules that preserve decidability
under a Type-3 producer deserves formal treatment (the ladder in
Section~\ref{sec:verifier-ladder} provides the analytical frame;
what is missing is a concrete domain that crosses the
Type-3/Type-2 boundary).

\subsection{Closing remark}

Turing wrote in 1936 that ``the justification [for the
definitions] lies in the fact that the human memory is
necessarily limited''~\cite{turing1936computable}.  Our laptop's
memory is necessarily limited.  Our agent's tape is necessarily
limited.  Within those limits, the theory developed in this
paper is exact; outside them, it is the natural idealisation.
An agent is its trace language --- nothing more, nothing less ---
and narrowness, paired with a producer--consumer discipline that
bounds the idealised Type-0 by the decidable Type-3, is enough.
